\section{Frontend}
\label{sec:frontend}
\textit{Responsable: Guillermo (FE-01, FE-02, FE-03).}

\pendiente{Describir el frontend (React + TypeScript + Vite en AWS Amplify): pantallas para docente y
estudiante, cómo se conecta a las APIs (variables \texttt{VITE\_API\_*} leídas de SSM por
\texttt{infra/frontend}) y el modo de autenticación (\texttt{VITE\_AUTH\_MODE}: mock o cognito).}

\subsection{Implementación y decisiones}
\pendiente{Estructura del proyecto, librerías, manejo de estados de carga y error, y cómo se pasa de
los mocks (MSW) a las APIs reales.}

\subsection{Resultados y evidencias}
\pendiente{URL pública de Amplify y una captura por pantalla principal del flujo del curso piloto.}

\evidencia{frontend-login.png}{Inicio de sesión con Amazon Cognito.}{fig:fe-login}
\evidencia{frontend-cursos.png}{Lista de cursos del docente.}{fig:fe-cursos}
\evidencia{frontend-modulo.png}{Módulo publicado con su material.}{fig:fe-modulo}
\evidencia{frontend-entrega.png}{Entrega de una tarea por el estudiante.}{fig:fe-entrega}
\evidencia{frontend-notas.png}{Notas del estudiante.}{fig:fe-notas}
